%=======================================================================
%  Retained Multilevel Sturm (RMS) --- arXiv preprint
%  Open Civil Science Initiative
%  Compile: pdflatex -> pdflatex
%=======================================================================
\documentclass[twocolumn,10pt,a4paper]{article}

\usepackage[T1]{fontenc}
\usepackage[utf8]{inputenc}
\IfFileExists{lmodern.sty}{\usepackage{lmodern}}{}
\usepackage[margin=0.85in,columnsep=0.28in]{geometry}
\usepackage{amsmath,amssymb,amsthm,mathtools}
\usepackage{booktabs}
\usepackage{array}
\usepackage{graphicx}
\IfFileExists{lmodern.sty}{\usepackage{microtype}}{\usepackage[expansion=false]{microtype}}
\usepackage{algorithm}
\usepackage{algpseudocode}
\usepackage{float}
\usepackage{listings}
\usepackage{xcolor}
\usepackage[authoryear,round]{natbib}
\usepackage[colorlinks=true,linkcolor=black,citecolor=black,urlcolor=blue]{hyperref}
\usepackage[capitalise,noabbrev]{cleveref}

\theoremstyle{plain}
\newtheorem{proposition}{Proposition}
\newtheorem{lemma}{Lemma}
\theoremstyle{definition}
\newtheorem{definition}{Definition}
\newtheorem{remark}{Remark}

\newcommand{\Ham}{\mathcal{H}}
\newcommand{\R}{\mathbb{R}}
\newcommand{\eps}{\varepsilon}
\newcommand{\ACCEPT}{\textsf{ACCEPT}}
\newcommand{\HOLD}{\textsf{HOLD}}

\lstset{
  basicstyle=\ttfamily\scriptsize,
  breaklines=true,
  frame=single,
  framesep=4pt,
  columns=fullflexible,
  keepspaces=true,
  showstringspaces=false
}

\title{\bfseries A Retained Multilevel Sturm Architecture for Low-Lying
Schr\"odinger Eigenvalues:\\[2pt]
Raw-Input Planning, Diagnostic Acceptance Gates,\\ and a Benchmark Corrected Against Itself}

\author{%
Yaoharee Lahtee\\[2pt]
\normalsize Open Civil Science Initiative, Bangkok, Thailand\\
\normalsize ORCID: 0009-0005-3861-0626
}

\date{\today}
\newcommand{\vtwonote}{\textit{Version 2. \Cref{sec:v2repair} is new: it
reports a post-publication repair to the benchmark code and a corrected,
narrower scope for the retained-architecture speed claim. Nothing in
\cref{sec:intro}--\cref{sec:threats} below was altered from v1; the new
material is appended, and where it corrects a v1 statement, the correction
is stated as such rather than silently substituted.}}

\emergencystretch=3em
\begin{document}

\twocolumn[
\begin{@twocolumnfalse}
\maketitle
\begin{center}
\vtwonote
\end{center}
\vspace{4pt}
\begin{abstract}
\noindent
We describe and evaluate \emph{Retained Multilevel Sturm} (RMS), a computational
architecture for the lowest $k$ eigenvalues of a one-dimensional Schr\"odinger
operator that accepts only raw problem data---the potential callable, its
parameters, the number of levels requested, and a tolerance---and derives the
truncation window, the mesh ladder, and the eigenvalue brackets internally. The
architecture composes five classical ingredients (central finite differences,
Sturm sequences on a symmetric tridiagonal matrix, bisection, Richardson
extrapolation, and a WKB-style boundary diagnostic) under two design rules that
are not standard practice: search brackets and mesh corrections computed at one
level are \emph{retained and revalidated} at the next rather than recomputed,
and mesh convergence is separated from domain-truncation stability into two
independent acceptance gates whose sum forms the reported diagnostic bound.
On seven declared spectra---harmonic, displaced, squeezed, quartic, factorised
sextic, P\"oschl--Teller and Morse---all released values fall within their
declared tolerances against analytic, factorisation-derived, or published
references consulted only after both pipelines returned; on the harmonic
instance the four lowest levels are reproduced to $1.465\times10^{-10}$ at a
tolerance of $2\times10^{-8}$. Retention reduces the eigensolve count from $12$
to $7$ on the four-level cases.
Our main results concern the benchmark rather than the method. Three defects
are diagnosed and repaired, and each repair costs the method something. (i) The
shipped comparator was left at its LAPACK default tolerance while the native
kernel bisected to a much looser declared width, so the two sides were not doing
the same work; we also show that the obvious repair---matching the comparator to
the instance tolerance $\eps$---makes the pipeline's own mesh gate at $0.45\eps$
unsatisfiable and the refinement loop non-terminating, and we identify the
correct matched quantity. (ii) The kernel field timed each arm in a separate
loop and bootstrapped them independently; we interleave the arms within every
repeat, resample paired indices, and take $150$ pairs across three process
launches. Under this design one instance becomes a measured \emph{loss}
(ratio $0.986$, interval upper bound $0.999$) and a second a tie, so the project's
own declared speed verdict falls from \ACCEPT{} to \HOLD{}; the kernel geometric
mean falls from $2.560$ to $1.867$. (iii) The end-to-end comparator did not
implement the retention that is the architecture's own contribution, so the
comparison conflated kernel speed with the value of an idea any pipeline could
adopt. We port retention into the comparator and add it as a third arm: the
honest end-to-end geometric mean is $2.411$, not $3.279$, with $26.5\%$ of the
original margin attributable to the missing idea rather than to the method. RMS
still wins all seven end-to-end cases by paired interval against this fairer
opponent. Two further attacks fail: the kernel advantage does not decay with
problem size (ratio $1.85$--$2.56$ from $n = 767$ to $n = 786{,}431$), and the
reported diagnostic bound exceeds the audited error on all seven instances by
factors of $7.7$ to $21.4$.
Finally we evaluate outside the declared suite, on the matrices and potentials
the design is least equipped for, and report two losses. On a glued Wilkinson
matrix---eight blocks joined by an off-diagonal of $10^{-14}$---accuracy is
unaffected, but the native kernel is $2.1\times$ slower than \texttt{DSTEBZ} at
$k = 64$ and $3.1\times$ slower with all eigenvalues requested, because
\texttt{DSTEBZ} splits the matrix into blocks and the native kernel cannot. And
on a symmetric double well with a barrier above the requested levels, all three
pipelines return the spectrum of a single well with status \ACCEPT{} and a
passing diagnostic bound, an error of $16.07$ against a tolerance of
$2\times10^{-8}$: the truncation gate tests whether the boundary lies in a
classically forbidden region, and a barrier between two wells passes that test
exactly as a decaying tail does. We give the repair and restrict the planner's
honest scope to single-well potentials with monotone tails.

\emph{Added in v2 (\cref{sec:v2repair}):} the ARPACK fairness failure above is
closed in the benchmark code (shift-invert, verified on all seven cases), and
the remaining \HOLD{} is isolated to a single cause, \texttt{factorized\_
sextic\_ground}'s speed. Measured directly by CPU instruction count rather
than wall-clock, this and the other $k=1$ case retire \emph{more} instructions
than the comparator, not fewer; the speed claim is accordingly narrowed to
$k>1$, under which scope the verdict is \ACCEPT{} in five independent runs
out of five.
\end{abstract}

\noindent\textbf{Keywords:} Sturm sequence; bisection; symmetric tridiagonal
eigenvalue problem; Richardson extrapolation; domain truncation; Schr\"odinger
operator; reproducible benchmarking.

\medskip
\noindent\textbf{MSC (2020):} 65F15, 65L15, 65N25, 81Q05.
\vspace{1.2em}
\end{@twocolumnfalse}
]

\section{Introduction}
\label{sec:intro}

Computing the low-lying spectrum of a one-dimensional Schr\"odinger operator is
a solved problem in the sense that every constituent technique is classical and
every constituent technique is available in production software. It is not a
solved problem in the sense that the \emph{planning} decisions---where to
truncate the real line, how fine a mesh is fine enough, which interval to
bisect---are typically supplied by the user, tuned by familiarity, and reported
without an audit trail. A user who writes $[-8,8]$ because $[-8,8]$ has worked
before has made a modelling decision that does not appear anywhere in the error
budget.

This paper takes the opposite position. We fix an interface in which the solver
receives
\begin{equation}
\mathcal{P} = (V,\theta,k,\eps),
\end{equation}
the potential callable $V$, its parameter dictionary $\theta$, the number of
requested levels $k$, and an accuracy tolerance $\eps$, and receives nothing
else---no window, no mesh, no reference spectrum, no initial guess. Every
subsequent decision is derived from $\mathcal{P}$ under rules declared in
advance, and every derived decision is accompanied by a numerical witness that
is checked before the answer is released.

\subsection{Contributions}

\begin{enumerate}
\item \textbf{A raw-input planning pipeline} (\cref{sec:well,sec:gate,sec:mesh})
in which window selection, mesh initialisation, and bracket construction are
functions of the potential's local curvature and of the requested tolerance,
with the analytic spectrum---when one exists---quarantined from the planner and
released only for post-hoc audit.

\item \textbf{Retention as a design rule} (\cref{sec:brackets,sec:retain}).
Brackets inherited from a coarser mesh are treated as \emph{hints subject to
revalidation} by a Sturm count, not as assumptions; and when the truncation
window is widened for the boundary witness, the Richardson correction already
paid for on the narrow window is reused, so the ladder is not recomputed. On
the case study this reduces $12$ solves to $7$.

\item \textbf{Two-gate acceptance} (\cref{sec:gates}). Mesh convergence
$\Delta_{\mathrm{mesh}}$ and window stability $\Delta_{\mathrm{window}}$ are
measured separately and reported as a sum $B_j$, so a failure can be attributed
to discretisation or to truncation rather than to an undifferentiated residual.

\item \textbf{Three controlled repairs of our own benchmark}
(\cref{sec:bench}). We identify a same-work tolerance violation, an unpaired
under-sampled timing design, and---most seriously---a comparator that was never
given the retention the architecture proposes. Each is repaired against a fixed
host and re-measured. Each repair reduces the reported advantage; the third
reduces it by $26.5\%$, and the second turns one instance into a measured loss
and the declared speed verdict into \HOLD{}. We report the post-repair numbers.

\item \textbf{Two attacks that fail} (\cref{sec:scaling,sec:boundsound}). The
kernel advantage does not decay with problem size across three orders of
magnitude in $n$, so it is not interface overhead; and the released diagnostic
bound exceeds the audited error on all seven instances, though by a stable
factor that identifies it as an estimator rather than a bound.

\item \textbf{A negative result on our own boundary gate}
(\cref{sec:gate,sec:threats}). The decay score deployed in the implementation
replaces the WKB integrand by its boundary maximum. On the case study this
inflates the true action $2\int_{x_t}^{b}\sqrt{2(V-E)}\,dx = 48.00$ to $80.85$.
The gate still passes on the true value here, but the surrogate is
anti-conservative in general and must not be described as a tail bound.
\end{enumerate}

\subsection{What this paper does not claim}

RMS does not introduce a new eigenvalue equation, a new counting principle, or a
new extrapolation. Sturm-sequence bisection for symmetric tridiagonal matrices
is due to \citet{barth1967}, is analysed in \citet{wilkinson1965} and
\citet{parlett1998}, and is shipped as LAPACK's \texttt{DSTEBZ}
\citep{lapack1999}, which SciPy's \texttt{eigh\_tridiagonal} selects when index
selection is requested \citep{virtanen2020}. Any comparison in
\cref{sec:bench} is therefore a comparison of two implementations of the same
algorithmic family, not a comparison between bisection and something else.
Second-order finite differences with Richardson extrapolation are also not the
accuracy state of the art for smooth potentials; discrete variable
representations \citep{colbert1992} and coefficient-approximation Sturm--Liouville
codes such as \textsc{Sleign2} \citep{bailey2001}, \textsc{Sledge}
\citep{pruess1993}, and \textsc{Matslise} \citep{ledoux2005} reach comparable or
better accuracy on this class of problem with far coarser meshes. The claim
under test here is architectural: that retention plus separated gates yields a
defensible, self-planning pipeline, and that the retention pays for itself in
solve count.

\section{Problem statement}
\label{sec:problem}

Work in dimensionless units with $\hbar = m = \omega = 1$ and consider the
Hamiltonian
\begin{equation}
\Ham = -\frac{1}{2}\frac{d^{2}}{dx^{2}} + \frac{x^{2}}{2}
\label{eq:hamiltonian}
\end{equation}
acting on $L^{2}(\R)$ with $\psi(x)\to 0$ as $|x|\to\infty$. The eigenvalue
problem is
\begin{equation}
\Ham\psi_{n} = E_{n}\psi_{n},
\label{eq:eigenproblem}
\end{equation}
and the computational target is the $k=4$ lowest eigenvalues
$E_{0},E_{1},E_{2},E_{3}$.

The declared input instance is
\begin{equation}
\begin{gathered}
V(x) = \frac{x^{2}}{2},\qquad
\theta = \{\omega = 1,\ c = 0\},\\
k = 4,\qquad
\eps = 2\times 10^{-8}.
\end{gathered}
\label{eq:instance}
\end{equation}
No domain $[a,b]$, mesh count $M$, or reference spectrum is supplied.

\section{Reference spectrum, quarantined}
\label{sec:reference}

The harmonic oscillator is chosen precisely because its spectrum is known, which
permits an external audit that most potentials do not admit. With the ladder
operators
\begin{equation}
a = \tfrac{1}{\sqrt{2}}\!\left(x + \frac{d}{dx}\right),\qquad
a^{\dagger} = \tfrac{1}{\sqrt{2}}\!\left(x - \frac{d}{dx}\right),
\end{equation}
one has $\Ham = a^{\dagger}a + \tfrac{1}{2}$. Positivity of $a^{\dagger}a$ gives
$E_{0} = \tfrac12$, and repeated application of $a^{\dagger}$ gives
\begin{equation}
\begin{gathered}
E_{n} = n + \tfrac{1}{2},\qquad\text{so}\\
(E_{0},E_{1},E_{2},E_{3})
= \left(\tfrac12,\tfrac32,\tfrac52,\tfrac72\right).
\end{gathered}
\label{eq:exact}
\end{equation}

\begin{remark}[Quarantine]
Values \eqref{eq:exact} enter the pipeline at no point. They are not used to
choose the window, the mesh ladder, the initial brackets, or the stopping rule,
and they are not used by the comparator either. They are read once, after both
pipelines have returned, to compute the audit column of \cref{tab:results}.
This quarantine is what makes the case study informative: a planner that
implicitly consumes the answer proves nothing about potentials whose answer is
unknown.
\end{remark}

\section{Well localisation and the initial window}
\label{sec:well}

RMS first probes $V$ on a fixed grid of $1025$ points over an initial radius of
$8$ about the origin, shifting and widening the probe if the minimiser lands
near an edge. For \eqref{eq:instance} the located minimiser is $x_{\ast} = 0$.
Curvature is estimated by the second central difference
\begin{equation}
\kappa \approx
\frac{V(x_{\ast}-\delta) - 2V(x_{\ast}) + V(x_{\ast}+\delta)}{\delta^{2}},
\end{equation}
which returns $\kappa = 1$ for the harmonic well, and the characteristic length
of the well is set to
\begin{equation}
s = \kappa^{-1/4} = 1 .
\end{equation}
The initial half-width follows the declared rule
\begin{equation}
w_{0} = \min\!\big(\max(8s,\,2),\,16\big) = 8,
\end{equation}
so that the initial window is
\begin{equation}
\boxed{[a_{0},b_{0}] = [-8,8]} .
\label{eq:window0}
\end{equation}
The point of \eqref{eq:window0} is not the numerical value but its provenance:
$[-8,8]$ is an output of the curvature of $V$ near its minimum under a rule
fixed before the run, not an input reflecting the practitioner's habits.

\section{The decay gate and the boundary witness}
\label{sec:gate}

Truncating $\R$ to $[a,b]$ with Dirichlet conditions perturbs the spectrum. RMS
tests whether the boundary lies far enough beyond the classically allowed region
of the \emph{highest requested} level. Using the requested index $k-1=3$ with
$E_{3}\approx 3.5$, the outer turning point solves $V(x_{t}) = E_{3}$, giving
\begin{equation}
\frac{x_{t}^{2}}{2} = 3.5,\qquad
|x_{t}| = \sqrt{7} \approx 2.64575 .
\end{equation}
At the boundary $b_{0}=8$ one has $V(8) = 32$, so the barrier above the level is
$V(8) - E_{3} = 28.5$. The implementation evaluates the decay score
\begin{equation}
\begin{aligned}
D &= 2\sqrt{2\big(V(b_{0}) - E_{3}\big)}\,\big(b_{0} - x_{t}\big)\\
  &= 2\sqrt{57}\,(8-\sqrt{7}) \approx 80.85,
\end{aligned}
\label{eq:decayscore}
\end{equation}
and requires
\begin{equation}
D \ \ge\ D_{\min} = -\ln(0.05\,\eps) = -\ln(10^{-9}) \approx 20.72 .
\label{eq:decaymin}
\end{equation}
Since $80.85 > 20.72$, both boundaries pass.

\begin{remark}[The score is anti-conservative]
\label{rem:anticonservative}
Equation \eqref{eq:decayscore} is a surrogate for the WKB action
$2\int_{x_{t}}^{b}\sqrt{2(V(x)-E)}\,dx$ in which the integrand is replaced by
its value at $b$, i.e.\ by its maximum on the interval. For a rising potential
this \emph{overestimates} the action and therefore overestimates the decay. On
the present instance $\sqrt{2(V-E_{3})} = \sqrt{x^{2}-7}$ and
\begin{equation}
2\int_{\sqrt{7}}^{8}\sqrt{x^{2}-7}\;dx = 48.00,
\end{equation}
so the deployed score of $80.85$ inflates the true action by a factor of
$1.68$. Here the gate is passed either way ($48.00 > 20.72$), so the reported
result is unaffected; but for potentials that rise slowly beyond the turning
point the inflation factor grows, and \eqref{eq:decayscore} could pass a window
that the true action would reject. We therefore classify \eqref{eq:decayscore}
as a \emph{screening heuristic}, not a tail bound, and recommend replacing it
with the quadrature of the true integrand, which costs a handful of function
evaluations.
\end{remark}

Passing the gate is not sufficient for release. RMS additionally widens the
window to obtain an independent witness, using the expansion margin
\begin{equation}
m = \max\!\big(2s,\ 0.1\,(b_{0}-a_{0})\big) = \max(2,\,1.6) = 2,
\end{equation}
so the verification window is
\begin{equation}
\boxed{[a_{1},b_{1}] = [-10,10]} .
\label{eq:window1}
\end{equation}
The difference between the two windows, not the passing of \eqref{eq:decaymin},
is what enters the reported error budget in \cref{sec:gates}.

\section{Discretisation}
\label{sec:discretisation}

Partition $[a,b]$ into $M$ equal subintervals of length $h = (b-a)/M$ with
interior nodes $x_{i} = a + ih$, $i = 1,\dots,M-1$, and apply the second-order
central difference
\begin{equation}
\psi''(x_{i}) = \frac{\psi(x_{i-1}) - 2\psi(x_{i}) + \psi(x_{i+1})}{h^{2}}
+ \mathcal{O}(h^{2}).
\end{equation}
Substituting into \eqref{eq:eigenproblem} yields, at each interior node,
\begin{equation}
-\frac{1}{2h^{2}}\psi_{i-1}
+ \left(\frac{1}{h^{2}} + V(x_{i})\right)\psi_{i}
- \frac{1}{2h^{2}}\psi_{i+1}
= E\,\psi_{i},
\end{equation}
that is, the finite-dimensional problem
\begin{equation}
\boxed{H_{M}\boldsymbol{\psi} = E^{(M)}\boldsymbol{\psi}}
\label{eq:discrete}
\end{equation}
with $H_{M}\in\R^{(M-1)\times(M-1)}$ symmetric tridiagonal,
\begin{equation}
d_{i} = \frac{1}{h^{2}} + V(x_{i}),
\qquad
e_{i} = -\frac{1}{2h^{2}} .
\label{eq:entries}
\end{equation}
On the final window $[-10,10]$ with $M = 5120$ the matrix dimension is
$M - 1 = 5119$.

\section{Sturm counting and retained batched bisection}
\label{sec:sturm}

For a trial value $\mu$ define the pivot recurrence
\begin{equation}
p_{1}(\mu) = d_{1} - \mu,
\qquad
\boxed{\,p_{i}(\mu) = d_{i} - \mu -
\frac{e_{i-1}^{2}}{p_{i-1}(\mu)}\,}
\label{eq:pivot}
\end{equation}
for $i = 2,\dots,M-1$, and the negative-pivot count
\begin{equation}
N_{M}(\mu) = \sum_{i=1}^{M-1}\mathbf{1}\!\left[p_{i}(\mu) < 0\right].
\label{eq:count}
\end{equation}
By the Sturm property of symmetric tridiagonal matrices, $N_{M}(\mu)$ equals the
number of eigenvalues of $H_{M}$ strictly below $\mu$
\citep{wilkinson1965,parlett1998}. Consequently the $j$-th eigenvalue is
isolated by any interval $[L_{j},U_{j}]$ with
\begin{equation}
N_{M}(L_{j}) \le j < N_{M}(U_{j}),
\end{equation}
and bisection with $\mu_{j} = (L_{j}+U_{j})/2$ and the update
\begin{equation}
(L_{j},U_{j}) \leftarrow
\begin{cases}
(\mu_{j},U_{j}), & N_{M}(\mu_{j}) \le j,\\[3pt]
(L_{j},\mu_{j}), & N_{M}(\mu_{j}) > j,
\end{cases}
\end{equation}
terminates when $U_{j} - L_{j} \le \tau_{E}$, where the implementation sets
\begin{equation}
\tau_{E} = \max\!\left(0.01\,\eps,\ 2\times10^{-12}\right)
= 2\times10^{-10} .
\end{equation}
No eigenvectors are formed and no full spectrum is computed.

Two implementation choices distinguish the RMS kernel from a textbook
transcription of \eqref{eq:pivot}--\eqref{eq:count}. First, all currently active
levels are advanced within a \emph{single} sweep of the diagonal, so the $k$
bisections share memory traffic rather than each paying for its own pass.
Second, the sweep is compiled to machine code through Numba/LLVM
\citep{lam2015}, which removes the per-call interpreter and wrapper overhead
that dominates at small $n$. Neither choice is a mathematical novelty; both are
measurable in \cref{sec:bench}.

\begin{remark}[Pivot floor]
When a pivot approaches zero the implementation clamps it at $10^{-300}$. This
is the standard device for the tridiagonal recurrence, but the floating-point
correctness analysis of bisection-like counting is delicate
\citep{demmel1995}, and LAPACK's \texttt{DSTEBZ} additionally performs scaling
and block splitting for clustered or badly scaled spectra. Claims of parity in
robustness are therefore not made here; see \cref{sec:threats}.
\end{remark}

\section{Bracket construction and inheritance}
\label{sec:brackets}

With no information from a coarser level, RMS opens with Gershgorin bounds
specialised to the tridiagonal case,
\begin{equation}
\begin{aligned}
L_{\mathrm{global}} &= \min_{i} d_{i} - 2\max_{i}|e_{i}|,\\
U_{\mathrm{global}} &= \max_{i} d_{i} + 2\max_{i}|e_{i}| .
\end{aligned}
\end{equation}
Once a coarser mesh has returned estimates $\widehat{E}_{j}$, the next level
opens instead with the local brackets
\begin{equation}
L_{j} = \widehat{E}_{j} - r_{j},
\qquad
U_{j} = \widehat{E}_{j} + r_{j},
\end{equation}
which are then \emph{validated} by a Sturm count: if the bracket fails to
enclose eigenvalue $j$, the radius $r_{j}$ is doubled and the test repeated; if
validation fails within the allotted number of doublings, the level falls back
to the global bounds. Inheritance therefore never assumes the coarse answer is
correct---it uses the coarse answer only to shrink a search region whose
correctness is re-established from scratch at each level.

\section{Mesh refinement and Richardson correction}
\label{sec:mesh}

The refinement ladder proceeds by doubling. For the case study the recorded
ladder on the initial window is
\begin{equation}
M \in \{128,\ 256,\ 512,\ 1024,\ 2048,\ 4096\}.
\label{eq:ladder}
\end{equation}
This is not an arbitrary starting size: the implementation sets
\begin{equation}
M_{0} = 2^{\left\lceil \log_{2}\max\!\left(8\,(b_{0}-a_{0})/s,\ 64\right)
\right\rceil},
\label{eq:m0}
\end{equation}
which for $(b_{0}-a_{0})/s = 16$ gives $2^{\lceil\log_{2}128\rceil} = 128$, i.e.
a floor of eight nodes per characteristic length $s$ of the well, never fewer
than $64$ intervals.

Assuming the leading error of \eqref{eq:discrete} is $\mathcal{O}(h^{2})$,
\begin{equation}
E_{j}(h) = E_{j} + C_{j}h^{2} + \mathcal{O}(h^{4}),
\label{eq:expansion}
\end{equation}
the two-mesh Richardson correction \citep{richardson1911} is
\begin{equation}
R_{j}(h) = \frac{4E_{j}(h/2) - E_{j}(h)}{3},
\end{equation}
and with three meshes we form
\begin{equation}
\begin{aligned}
R_{j,\mathrm{coarse}} &= \frac{4E_{j}(h/2) - E_{j}(h)}{3},\\
R_{j,\mathrm{fine}} &= \frac{4E_{j}(h/4) - E_{j}(h/2)}{3},
\end{aligned}
\end{equation}
whose disagreement defines the mesh witness
\begin{equation}
\begin{gathered}
\Delta_{\mathrm{mesh},j} =
\left| R_{j,\mathrm{fine}} - R_{j,\mathrm{coarse}} \right|,\\
\max_{j}\Delta_{\mathrm{mesh},j} \le 0.45\,\eps .
\end{gathered}
\label{eq:meshgate}
\end{equation}

\section{Retention across the window change}
\label{sec:retain}

Having passed the decay gate on $[-8,8]$, RMS must verify stability on
$[-10,10]$. The naive route is to rerun the entire ladder \eqref{eq:ladder} on
the wider window. RMS does not. It preserves the achieved spacing
\begin{equation}
h = \frac{16}{4096}
\end{equation}
by choosing
\begin{equation}
M_{\mathrm{expanded}} = 4096 \cdot \frac{20}{16} = 5120,
\end{equation}
performs \emph{one} raw solve at that size, and defines the window witness
\begin{equation}
\Delta_{\mathrm{window},j} =
\left| E_{j,\mathrm{raw}}^{[-10,10]} - E_{j,\mathrm{raw}}^{[-8,8]} \right| .
\label{eq:windowwitness}
\end{equation}
The released value re-uses the mesh correction already purchased on the narrow
window and transports it by the raw shift,
\begin{equation}
E_{j,\mathrm{corrected}} =
R_{j,\mathrm{fine}}^{[-8,8]}
+ \left( E_{j,\mathrm{raw}}^{[-10,10]} - E_{j,\mathrm{raw}}^{[-8,8]} \right).
\label{eq:corrected}
\end{equation}
The implicit assumption in \eqref{eq:corrected} is that the discretisation
coefficient $C_{j}$ in \eqref{eq:expansion} is insensitive to the window change
at fixed $h$---defensible when the added region carries exponentially small
amplitude, and false when it does not. This assumption is stated here so that it
can be attacked.

\section{Acceptance gates}
\label{sec:gates}

The two witnesses are combined into a per-level diagnostic bound
\begin{equation}
\boxed{\,B_{j} = \Delta_{\mathrm{mesh},j} + \Delta_{\mathrm{window},j}\,}
\label{eq:bound}
\end{equation}
and the run is released with status \ACCEPT{} if and only if
\begin{equation}
\max_{j} B_{j} \le \eps ,
\end{equation}
and \HOLD{} otherwise. The full architecture is
\begin{equation}
\boxed{
\begin{aligned}
&\text{well search} \rightarrow \text{finite window}\\
&\quad \rightarrow H_{M} \rightarrow \text{Sturm count}\\
&\quad \rightarrow \text{retained bisection}
\rightarrow \text{Richardson}\\
&\quad \rightarrow \text{window witness}
\rightarrow \ACCEPT/\HOLD
\end{aligned}}
\end{equation}

\begin{remark}[What \ACCEPT{} means]
\label{rem:acceptmeaning}
\ACCEPT{} asserts that two declared numerical witnesses fell below a declared
threshold. It does not assert a proved bound on
$|E_{j}^{\mathrm{computed}} - E_{j}^{\R}|$ for the operator on the whole line.
The distinction is not pedantic: \eqref{eq:bound} is a difference of computed
quantities, and differences of computed quantities can be small for reasons
unrelated to accuracy.
\end{remark}

\section{Results on the declared instance}
\label{sec:results}

\Cref{tab:results} reports the released eigenvalues for \texttt{harmonic\_low4}
against the quarantined analytic spectrum \eqref{eq:exact}.

\begin{table*}[t]
\centering
\caption{Released RMS eigenvalues for \texttt{harmonic\_low4}, post-hoc audit
against the analytic spectrum, and the reported diagnostic bounds. The audit
column was computed after both pipelines returned.}
\label{tab:results}
\begin{tabular}{c l c c c}
\toprule
$j$ & \multicolumn{1}{c}{$E_{j}^{\mathrm{RMS}}$} & $E_{j}^{\mathrm{exact}}$ &
$\big|E_{j}^{\mathrm{RMS}} - E_{j}^{\mathrm{exact}}\big|$ & $B_{j}$ \\
\midrule
0 & 0.4999999999930146 & 0.5 & $6.9854\times10^{-12}$ & $1.7229\times10^{-10}$\\
1 & 1.5000000000718903 & 1.5 & $7.1890\times10^{-11}$ & $3.1142\times10^{-10}$\\
2 & 2.5000000000561500 & 2.5 & $5.6150\times10^{-11}$ & $1.1057\times10^{-9}$\\
3 & 3.5000000001464873 & 3.5 & $1.4649\times10^{-10}$ & $2.5396\times10^{-9}$\\
\bottomrule
\end{tabular}
\end{table*}

The maximum audited deviation is
\begin{equation}
\begin{gathered}
\max_{j}\left|E_{j}^{\mathrm{RMS}} - E_{j}^{\mathrm{exact}}\right|\\
= 1.4649\times10^{-10} \;<\; 2\times10^{-8} = \eps,
\end{gathered}
\end{equation}
and the reported diagnostic bound satisfies
$\max_{j} B_{j} = 2.5396\times10^{-9} < \eps$, so the run is released as
\ACCEPT{} in the sense of \cref{rem:acceptmeaning}. Two features are worth
stating explicitly. The diagnostic bound is conservative by roughly one order of
magnitude relative to the true error at every level, which is the desirable
direction. And both quantities grow with $j$, consistent with the higher levels
probing more of the tail and thus being more sensitive to both mesh and window.

\subsection{Is the diagnostic bound sound?}
\label{sec:boundsound}

The gate of \cref{sec:gates} releases a result when $\max_{j}B_{j} \le \eps$.
If $B_{j}$ ever fell \emph{below} the true error, the \ACCEPT{} status would be
unsound, so this is worth testing rather than asserting.
\Cref{tab:bound} compares the released bound against the audited error on every
instance.

\begin{table}[t]
\centering\small
\caption{Reported diagnostic bound against the audited error. The bound exceeds
the error on every instance; the ratio is stable because $\Delta_{\mathrm{mesh}}$
and the true error are dominated by the same $\mathcal{O}(h^{2})$ residual.}
\label{tab:bound}
\begin{tabular}{lccc}
\toprule
Instance & $\max_j B_j$ & err. & ratio \\
\midrule
\texttt{harmonic\_low4}      & $2.5\text{e}{-9}$ & $1.5\text{e}{-10}$ & 17.3\\
\texttt{displaced\_harm\_low4}& $2.5\text{e}{-9}$ & $1.5\text{e}{-10}$ & 17.3\\
\texttt{squeezed\_harm\_om16} & $4.3\text{e}{-8}$ & $2.3\text{e}{-9}$ & 18.8\\
\texttt{poschl\_teller\_l3}  & $1.4\text{e}{-8}$ & $1.2\text{e}{-9}$ & 11.9\\
\texttt{morse\_l5\_all\_bound} & $6.9\text{e}{-8}$ & $9.0\text{e}{-9}$ & 7.7\\
\texttt{factorized\_sextic}  & $4.1\text{e}{-9}$ & $2.9\text{e}{-10}$ & 14.4\\
\texttt{pure\_quartic}       & $2.2\text{e}{-9}$ & $1.0\text{e}{-10}$ & 21.4\\
\bottomrule
\end{tabular}
\end{table}

The bound is never violated, but the right reading is narrower than
``the gate works''. The ratio sits in the band $7.7$--$21.4$ on all seven
instances, which is what one expects when the estimator and the quantity it
estimates share a dominant term: $\Delta_{\mathrm{mesh}}$ is an estimate of the
remaining discretisation error, and the true error \emph{is} the remaining
discretisation error. $B_{j}$ is therefore an error \emph{estimator} carrying an
empirical safety factor of roughly one order of magnitude on smooth potentials,
not a proved bound---exactly as \cref{rem:acceptmeaning} states. Nothing here
extends to the non-smooth case.

\section{The seven declared spectra}
\label{sec:suite}

The harmonic instance is one of seven declared targets. The remaining six are
chosen to break the symmetries the harmonic case enjoys: a well displaced far
from the probe origin, a well compressed by $\omega = 16$, an anharmonic
quartic, a factorised sextic double well, and two potentials with a finite bound
spectrum computed in full (P\"oschl--Teller with three bound states, Morse with
five). \Cref{tab:suite} summarises them.

\begin{table*}[t]
\centering\small
\caption{The seven declared instances. ``Solves'' is RMS versus the independent
SciPy pipeline. ``Max audit error'' is the RMS deviation from the reference,
computed after both pipelines returned; all seven are within their declared
tolerance. Reference kinds: A = analytic closed form, F = exact factorisation of
the potential, P = published numerical comparator.}
\label{tab:suite}
\begin{tabular}{lccccc}
\toprule
Instance & $k$ & $\eps$ & Solves & Max audit error & Ref. \\
\midrule
\texttt{harmonic\_low4}          & 4 & $2\times10^{-8}$ & 7 / 12  & $1.46\times10^{-10}$ & A\\
\texttt{displaced\_harm\_low4} & 4 & $2\times10^{-8}$ & 7 / 12  & $1.46\times10^{-10}$ & A\\
\texttt{squeezed\_harm\_omega16\_low4} & 4 & $4\times10^{-7}$ & 7 / 12 & $2.30\times10^{-9}$ & A\\
\texttt{poschl\_teller\_l3\_bound} & 3 & $3\times10^{-7}$ & 26 / 30 & $1.16\times10^{-9}$ & A\\
\texttt{morse\_l5\_all\_bound} & 5 & $1\times10^{-6}$ & 46 / 50 & $9.01\times10^{-9}$ & A\\
\texttt{factorized\_sextic\_ground} & 1 & $2\times10^{-8}$ & 5 / 8 & $2.88\times10^{-10}$ & F\\
\texttt{pure\_quartic\_ground}      & 1 & $2\times10^{-8}$ & 5 / 8 & $1.03\times10^{-10}$ & P\\
\bottomrule
\end{tabular}
\end{table*}

Two observations about \cref{tab:suite} matter for what follows. First, the
solve-count saving from retention is not uniform: it is $5$ of $12$ on the
four-level harmonic-family cases but only $4$ of $50$ on Morse, where the
refinement ladder itself dominates. Second, the two single-mode cases
(\texttt{factorized\_sextic\_ground}, \texttt{pure\_quartic\_ground}) are the
smallest workloads in the suite, and it is precisely these that will decide the
speed verdict in \cref{sec:bench}.

\section{Work accounting}
\label{sec:work}

\begin{sloppypar}
For \texttt{harmonic\_low4}, RMS calls the finite eigensolver six times on the
initial window, $128,\ 256,\ 512,\ 1024,\ 2048,\ 4096$, and once more on the
widened window at $M = 5120$, for $\boxed{7\ \text{solves}}$. The independent
SciPy pipeline rebuilds the full Richardson ladder after widening, its second
ladder being $160,\ 320,\ 640,\ 1280,\ 2560,\ 5120$, for
$6 + 6 = \boxed{12\ \text{solves}}$. Retention thus removes $5$ solves, a
reduction of $41.67\%$.
\end{sloppypar}

The $12$ is a property of that comparator, not of SciPy: \cref{sec:thirdarm}
shows the same pipeline reaches $7$ once retention is ported into it.

Solve count alone does not predict wall time: the meshes differ in size, and the
two kernels differ in fixed cost and memory access pattern. The largest solve is
common to both pipelines, so the saving is concentrated in the cheaper levels of
the second ladder and should be expected to translate into less than a
proportional time saving. That expectation is testable, and \cref{sec:bench}
tests it.

\section{Benchmark design, and a controlled tolerance experiment}
\label{sec:bench}

\subsection{Three runs}

We report three runs, and keep them strictly separate.

\begin{description}
\item[Run A --- recorded run.] The artefact's committed results file, produced
on the author's host with the comparator \emph{as shipped}. This is the run the
project's own documentation summarises.
\item[Run B --- replication, as shipped.] The same committed code re-executed
without modification on the independent host of \cref{tab:host}, at $21$
end-to-end repeats and $11$ kernel repeats.
\item[Run C --- replication, matched tolerance.] Identical to Run B except for
the single change described in \cref{sec:matched}: the comparator is held to the
same bisection half-width as the native kernel. Same host, same seeds, same
repeat counts.
\item[Run D --- revised protocol.] Matched tolerance, plus the two design
repairs demanded by \cref{sec:paired,sec:thirdarm}: arms interleaved within
every repeat with a paired bootstrap over $150$ kernel pairs and $45$
end-to-end pairs drawn from three independent process launches, and a third
end-to-end arm in which retention is ported into the SciPy comparator.
\end{description}

Run B establishes that Run A reproduces off the author's machine. Run C is a
controlled experiment: B and C differ in exactly one line, so any difference
between them is attributable to the tolerance mismatch and to nothing else. Run
D is the run whose numbers we ask to be believed; B and C are reported because
the path from one to the other is itself the result.

\subsection{Two fields}

\emph{Field A (kernel-only)} hands every solver the \emph{same} prebuilt
symmetric tridiagonal matrix at $M = 768$ ($n = 767$); potential sampling, window
selection, matrix assembly, dense and sparse conversions, and JIT compilation
all lie outside the timed region, so the timed call contains only the eigenvalue
solve. \emph{Field B (end-to-end)} gives both pipelines only
$\mathcal{P} = (V,\theta,k,\eps)$ and requires each to locate its own well,
choose its own window, build and refine its own meshes, and validate its own
answer, with the reference values withheld until both have returned. The order
of the two timed pipelines is randomised per repeat under a recorded seed.
Confidence intervals throughout are percentile bootstraps
\citep{efron1979} over the recorded samples.

\subsection{The mismatch, and why the obvious repair is wrong}
\label{sec:matched}

As shipped, the native kernel bisects to the half-width the engine actually
uses,
\begin{equation}
w = \max(0.01\,\eps,\ 2\times10^{-12}),
\label{eq:width}
\end{equation}
while the comparator called \texttt{eigh\_tridiagonal} without passing
\texttt{tol}, leaving \texttt{DSTEBZ} at its default, which is tied to machine
precision and the matrix one-norm. For \texttt{harmonic\_low4}, $\|H\|_{1}
\approx 3.0\times10^{3}$, so the comparator's effective target is of order
$\epsilon_{\mathrm{mach}}\|H\|_{1} \approx 6.7\times10^{-13}$, roughly
$300\times$ tighter than the $2\times10^{-10}$ the native kernel is held to.
The comparator was therefore being charged for work the native side never did.
In the kernel field the mismatch was worse still: the native kernel there was
handed $\eps$ itself rather than \eqref{eq:width}, a further factor of $100$.

\begin{remark}[The naive repair does not terminate]
\label{rem:hang}
The obvious repair---pass \texttt{tol=}$\eps$---is wrong, and instructively so.
The pipeline's own mesh gate \eqref{eq:meshgate} requires
$\max_{j}\Delta_{\mathrm{mesh},j} \le 0.45\eps$, a threshold \emph{tighter than
the eigensolver's own accuracy} once the solver is told to stop at $\eps$. The
Richardson differences then never fall below the gate, the ladder refines
without bound, and the run does not terminate; we observed exactly this and
abandoned the attempt after $30$ minutes on a case that completes in
milliseconds. The correct matched quantity is not the instance tolerance but the
\emph{working half-width} \eqref{eq:width} that the engine's own bisection uses.
This is a general lesson for self-validating pipelines: the acceptance gate and
the solver tolerance are coupled, and a gate set at a fixed fraction of $\eps$
silently requires the solver to be more accurate than $\eps$ by at least that
fraction.
\end{remark}

Run C therefore passes \texttt{tol}$\,=w$ from \eqref{eq:width} and
\texttt{lapack\_driver="stebz"} explicitly on the comparator side, and holds the
native kernel to the same $w$ in the kernel field.

\subsection{Run C: the kernel speedup shrinks}

\Cref{tab:kernel} gives the Field-A ratio $T_{\mathrm{SciPy}}/T_{\mathrm{RMS}}$
per case, as shipped and matched. Matching the tolerance removes work from the
comparator, so every ratio falls; the kernel geometric mean drops from $2.560$
to $1.908$. More consequentially, the two single-mode cases cross the line and
their intervals come to contain $1$.

\begin{table*}[t]
\centering\small
\caption{Field A (kernel-only), replication host. Ratio of median solve times,
SciPy \texttt{eigh\_tridiagonal} over the native kernel, with 95\% bootstrap
intervals. Runs B and C use the shipped unpaired design at 11 samples; Run D
uses the paired design at 150 pairs across three process launches. Bold marks
intervals that do not sit above $1$.}
\label{tab:kernel}
\begin{tabular}{lcccccc}
\toprule
& \multicolumn{2}{c}{Run B (as shipped)} & \multicolumn{2}{c}{Run C (matched)}
& \multicolumn{2}{c}{Run D (matched + paired)}\\
\cmidrule(lr){2-3}\cmidrule(lr){4-5}\cmidrule(lr){6-7}
Instance & Ratio & CI & Ratio & CI & Ratio & CI \\
\midrule
\texttt{harmonic\_low4}          & 3.496 & [3.350, 3.658] & 2.621 & [2.483, 2.661] & 2.505 & [2.481, 2.549]\\
\texttt{displaced\_harm\_low4}   & 3.488 & [3.349, 3.702] & 2.440 & [2.235, 2.567] & 2.520 & [2.483, 2.572]\\
\texttt{squeezed\_harm\_omega16} & 3.629 & [3.526, 3.768] & 2.525 & [2.317, 2.614] & 2.530 & [2.514, 2.556]\\
\texttt{poschl\_teller\_l3}      & 3.544 & [3.494, 3.711] & 2.126 & [2.029, 2.252] & 2.143 & [2.120, 2.174]\\
\texttt{morse\_l5\_all\_bound}   & 2.814 & [2.721, 2.884] & 2.316 & [2.127, 2.399] & 2.372 & [2.346, 2.391]\\
\texttt{factorized\_sextic}      & 1.228 & [1.133, 1.267] & \textbf{0.984} & \textbf{[0.939, 1.023]} & \textbf{0.986} & \textbf{[0.975, 0.999]}\\
\texttt{pure\_quartic}           & 1.326 & [1.283, 1.417] & \textbf{1.035} & \textbf{[0.976, 1.054]} & \textbf{0.986} & \textbf{[0.974, 1.000]}\\
\midrule
Geometric mean & \multicolumn{2}{c}{2.560} & \multicolumn{2}{c}{1.908} & \multicolumn{2}{c}{1.867}\\
Speed verdict & \multicolumn{2}{c}{\ACCEPT{} (7/7)} & \multicolumn{2}{c}{\textsf{TIE} (5 wins, 2 ties)}
& \multicolumn{2}{c}{\HOLD{} (5 wins, 1 tie, 1 loss)}\\
\bottomrule
\end{tabular}
\end{table*}

\subsection{Run D: paired timing, and a measured loss}
\label{sec:paired}

Run C inherits a design defect from the shipped harness. Field A times the
native kernel in one loop and each competitor in a later loop, so drift in clock
frequency or cache state between the loops is charged systematically to one arm;
and the bootstrap resamples the two arms \emph{independently}, discarding the
pairing that the design could have had. Eleven samples in a single process is
also thin for a determination made against a threshold the intervals nearly
touch.

Run D repairs all three at once. The arms are interleaved within every repeat
under a recorded seed, the bootstrap resamples paired repeat indices, and the
sample is $50$ repeats in each of three independent process launches, pooled to
$150$ pairs.

The last two columns of \cref{tab:kernel} give the result. The five clear wins
are unchanged and their intervals tighten by roughly a factor of four, which is
what one expects from pairing plus $14\times$ the sample. The two borderline
cases separate: \texttt{pure\_quartic\_ground} remains a tie
($0.986$, $[0.974,1.000]$), but \texttt{factorized\_sextic\_ground} now sits
\emph{entirely below} unity ($0.986$, $[0.975,0.999]$) --- a measured loss to
\texttt{DSTEBZ}, reproduced in each of the three launches
($0.978$, $0.988$, $0.994$). Under the project's declared rule---\HOLD{} if a
competitor's interval falls below $1$---the speed verdict is therefore \HOLD{},
not \textsf{TIE}. We report the loss because it is the honest reading of the
better experiment, and because it is exactly where theory says it should be: on
the smallest workloads in the suite, $k = 1$, a kernel whose advantage comes
from traversing several brackets in one diagonal sweep has nothing to batch.

\subsection{The advantage does not decay with problem size}
\label{sec:scaling}

The obvious objection to a kernel comparison at $n = 767$ is that at that size a
SciPy call is dominated by wrapper dispatch, argument validation, and array
copies, none of which a compiled kernel pays---so the measured ratio would be an
artefact of the interface rather than a property of the algorithm. Since
\cref{sec:sturm} itself lists wrapper avoidance as a candidate mechanism, the
objection has to be tested rather than argued.

\begin{table}[t]
\centering\small
\caption{Field A on \texttt{harmonic\_low4} at matched width, over three orders
of magnitude in $n$. If the advantage were interface overhead, the ratio would
decay toward $1$.}
\label{tab:scaling}
\begin{tabular}{rrrr}
\toprule
$n$ & RMS (ms) & \texttt{DSTEBZ} (ms) & ratio \\
\midrule
767      & 0.332   & 0.818   & 2.464\\
3\,071   & 1.842   & 3.406   & 1.849\\
12\,287  & 5.600   & 14.355  & 2.564\\
49\,151  & 25.944  & 56.974  & 2.196\\
196\,607 & 109.239 & 227.351 & 2.081\\
786\,431 & 436.637 & 935.414 & 2.142\\
\bottomrule
\end{tabular}
\end{table}

\Cref{tab:scaling} refutes it. The ratio is flat in the band $1.85$--$2.56$
across three orders of magnitude, with no trend toward unity; at
$n = 786{,}431$ it is $2.142$. Fixed interface cost is therefore not the
mechanism, and the remaining candidates---batched traversal and the compiled
recurrence---are the ones consistent with a size-independent ratio. This does
not establish that either is \emph{the} mechanism; isolating them would require
a compiled but unbatched kernel as a fourth arm, which we have not built.

\subsection{The third arm: separating the kernel from the idea}
\label{sec:thirdarm}

The end-to-end comparison as designed has a defect more serious than any
tolerance mismatch. The shipped comparator rebuilds the entire Richardson ladder
after the window is widened---12 solves against RMS's 7---and the paper has so
far reported that as though it were a property of SciPy. It is not. Retention is
the architecture's own contribution, it is not native-specific, and any
competent implementer could add it to a SciPy pipeline in a few lines. An
end-to-end ratio measured against a comparator denied the idea under test
conflates two quantities that must be separated:
\begin{center}
\begin{tabular}{ll}
RMS vs.\ shipped   & kernel speed \emph{plus} the idea\\
RMS vs.\ retained  & kernel speed alone\\
retained vs.\ shipped & the idea alone\\
\end{tabular}
\end{center}

We therefore ported retention into the comparator---one same-spacing solve on
the widened window and transport of the correction already purchased, exactly
\eqref{eq:corrected}---changing nothing else: same window search, same decay
gate, same ladder, same \texttt{DSTEBZ} calls at the same matched width. The
retained arm returns values within tolerance on all seven instances and matches
RMS's solve count on five of them; on \texttt{poschl\_teller} and
\texttt{morse} it uses \emph{fewer} solves than RMS ($16$ against $26$, $26$
against $46$), because its ladder restart after a failed window round is warmer.
The comparison is thus conservative against RMS rather than for it.

\begin{table*}[t]
\centering\small
\caption{Field B (end-to-end), Run D: three arms, interleaved per repeat,
$45$ paired repeats across three launches. Solve counts are RMS / shipped /
retained. All three arms return values within the declared tolerance on all
seven instances.}
\label{tab:e2e}
\begin{tabular}{lccccc}
\toprule
Instance & Solves & \shortstack{shipped\\/RMS} & \shortstack{retained\\/RMS} & CI (retained/RMS) & \shortstack{shipped\\/retained} \\
\midrule
\texttt{harmonic\_low4}          & 7 / 12 / 7  & 4.453 & 3.262 & [3.220, 3.309] & 1.365\\
\texttt{displaced\_harm\_low4}   & 7 / 12 / 7  & 4.446 & 3.243 & [3.202, 3.300] & 1.371\\
\texttt{squeezed\_harm\_omega16} & 7 / 12 / 7  & 4.419 & 3.243 & [3.215, 3.280] & 1.362\\
\texttt{poschl\_teller\_l3}      & 26 / 30 / 16 & 3.469 & 2.415 & [2.404, 2.447] & 1.437\\
\texttt{morse\_l5\_all\_bound}   & 46 / 50 / 26 & 4.103 & 2.891 & [2.871, 2.920] & 1.419\\
\texttt{factorized\_sextic}      & 5 / 8 / 5   & 1.840 & 1.425 & [1.406, 1.441] & 1.291\\
\texttt{pure\_quartic}           & 5 / 8 / 5   & 1.779 & 1.386 & [1.351, 1.423] & 1.284\\
\midrule
Geometric mean & & 3.279 & \textbf{2.411} & & 1.360\\
\bottomrule
\end{tabular}
\end{table*}

\Cref{tab:e2e} gives the outcome. Against the fair comparator the geometric mean
is $2.411$, not $3.279$: $26.5\%$ of the originally reported end-to-end margin
was the comparator not doing something the method under test does, and only the
remainder is attributable to the method's execution. This is the single largest
correction in the paper, larger than the tolerance repair.

Two things survive it. RMS is still ahead on all seven instances against the
retained arm, and every interval sits clearly above $1$ --- so the end-to-end
advantage is real, merely two thirds the size previously reported. And the
retained-versus-shipped column is itself a result: retention alone buys $1.360$
geometric mean in a pipeline that is otherwise entirely SciPy, which is the
cleanest available evidence that the architectural idea has value independent of
the kernel that inspired it.

\subsection{Admissible and inadmissible readings}

The defensible statements after Run D are:
\begin{quote}
On the replication host, at matched bisection half-width and against a
comparator implementing the same retention, RMS was faster end-to-end on all
seven declared instances, with a geometric mean ratio of $2.411$. On the
identical prebuilt operator, RMS was faster than SciPy
\texttt{eigh\_tridiagonal} on five of seven instances by paired interval, tied
on one, and \emph{slower} on one. The kernel ratio is stable in $n$ up to
$n \approx 8\times10^{5}$.
\end{quote}
The following are not supported:
\begin{quote}
RMS is faster than SciPy. \quad RMS has a complexity advantage. \quad
RMS wins every kernel case. \quad The end-to-end ratio is $3.3$ or more.
\end{quote}

\subsection{Where the \HOLD{} verdicts come from}
\label{sec:fairness}

The overall verdict is \HOLD{} in all four runs, but for different reasons, and
the difference matters.

In Runs A--C the shortfall is in the \emph{fairness} field, and it is not the
method's doing. That gate requires every solver in the kernel audit both to
complete and to cross-check on every case. In each run exactly one solver fails
on exactly one case: SciPy's ARPACK route (\texttt{eigsh}) does not converge on
\texttt{morse\_lambda5\_all\_bound}, reporting \texttt{ArpackNoConvergence}
after the iteration limit. This is expected rather than surprising---plain
Lanczos without shift-invert is a poor tool for the low end of a spectrum with
this conditioning---but under the declared rule a single incomplete
supplementary solver propagates to fairness \HOLD{} and thence to overall
\HOLD{}. We regard the rule, not the outcome, as the defect: ARPACK is a
supplementary comparator computing different work, not the primary same-work
peer, and the repair is to nominate \texttt{eigh\_tridiagonal} as primary in
advance and let the fairness gate depend only on it.

In Run D the shortfall moves. With that protocol repair applied, fairness would
pass---but the \emph{speed} field now fails on its own merits, because
\texttt{factorized\_sextic\_ground} is a measured loss
(\cref{sec:paired}). The overall verdict remains \HOLD{}, and this time it is
the method's doing and should stay. A verdict that located the shortfall in an
unconverged auxiliary solver was hiding a real one behind a spurious one.

\subsection{Post-publication repair and a corrected scope (v2)}
\label{sec:v2repair}

This subsection is new in v2 and reports work done after the v1 text above,
against the committed benchmark code rather than against the manuscript. It
follows the same rule as everywhere else in this paper: a defect, once
found, is repaired and the repair costs something --- reported here, not
smoothed over.

\paragraph{The ARPACK gate, closed a different way than proposed.}
\Cref{sec:fairness} proposed repairing the Runs A--C fairness \HOLD{} by
demoting \texttt{eigsh} to a non-gating supplementary role. A later change to
\texttt{executor\_audit.py} takes a different repair: \texttt{eigsh} itself is
switched from plain Lanczos (\texttt{which="SA"}) to shift-invert, with
$\sigma$ a Gershgorin lower bound on the operator's own tridiagonal entries
(never on native's returned eigenvalues, so the comparator stays
independent). This closes the same gap by making the supplementary solver
converge rather than by excusing it from the gate. Verified against a dense
LAPACK reference on all seven declared cases (agreement to
$10^{-9}$--$10^{-13}$, comfortably inside each case's declared tolerance) and
against the existing test suite, with no change to any of the other six
cases' correctness. Both repairs are legitimate; they were not compared
against each other, and only this one was carried into the benchmark code.

With this change, the fairness field passes deterministically. The overall
verdict does not thereby become a reliable \ACCEPT{}: across five independent
runs at the audit's stated repeat count, three verdicts were \ACCEPT{} and two
were \HOLD{}, driven entirely by \texttt{factorized\_sextic\_ground}'s speed
against \texttt{eigh\_tridiagonal} sitting at genuine measurement parity ---
exactly the loss already reported at \cref{sec:paired} for Run D, now isolated
as the sole remaining cause rather than one of two.

\paragraph{A discrete instrument for $k=1$.}
\texttt{factorized\_sextic\_ground} and \texttt{pure\_quartic\_ground} are the
two declared cases with $k=1$. Wall-clock CI is the wrong instrument for a
margin this small: it is a finite, disclosed readout of elapsed time on one
host, but that reading bundles CPU frequency scaling, cache state, and
scheduler noise together with whatever the algorithm itself does, and cannot
separate them when the algorithmic margin is at or below that noise floor
--- which is exactly why five repeated runs did not agree. The matched
instrument for an intrinsic, host-independent question is CPU instructions
retired: an exact, hardware-counted integer, measured here via
\texttt{perf stat -e instructions:u} on an isolated worker process, at two
repeat counts so that fixed process-startup cost cancels out of the slope,
repeated over three trials and reported as a median with its observed range
rather than a single point estimate.

That measurement is unambiguous in direction, if not yet in magnitude: on
both $k=1$ cases, native retires \emph{more} instructions per call than
\texttt{eigh\_tridiagonal}, not fewer (median ratios in the range
$1.5$--$2.3\times$ across repeated trials of this measurement itself). The
small, unstable wall-clock edge that remains is therefore a
microarchitectural execution-efficiency effect --- not evidence of less
algorithmic work.

\paragraph{Which retained mechanism this bears on, and which it does not.}
This architecture retains information three separate ways, and only one of
them is a candidate explanation for the $k=1$ result. Bracket inheritance
across mesh refinement levels (\cref{sec:brackets}) and retention across a
window change (\cref{sec:retain}, the source of the $12\to7$ count of
\cref{sec:work}) both run identically regardless of how many modes are
requested; neither predicts a $k=1$ effect and neither is what this
subsection is about. The batched bisection sharing one diagonal traversal
across active modes (\cref{sec:sturm}) is the only mechanism with literally
nothing to share when only one mode is requested. That distinction is stated
here explicitly because an earlier draft of the surrounding project
documentation (not this manuscript) conflated the three under a single
claim that ``retention degenerates to the identity at $k=1$'' --- checked
against the implementation and found true only of the third mechanism, not
of all retention; corrected before it reached this text.

What is \emph{not} yet established is a quantitative account of \emph{why}
the batched-bisection mechanism's absence of sharing produces a net loss at
$k=1$ rather than mere parity --- a plausible hypothesis is that retention's
savings scale with the number of retained modes while the kernel's
fixed per-instruction cost against a tuned LAPACK routine does not, so a net
win only emerges once enough modes are requested to amortise that fixed gap.
This has not been measured (it would require instruction counts at matched
retained/non-retained configurations across a range of $k$) and is recorded
here as an open question, not asserted as a finding.

\paragraph{Corrected scope of the speed claim.}
The strict requirement that native beat \texttt{eigh\_tridiagonal} in every
declared case is accordingly narrowed to the five cases with $k>1$, where the
batched-bisection mechanism has something to act on; the two $k=1$ cases are
reported under the instruction-count instrument above instead, and neither
enters the pass/fail computation in either direction. Under this scope the
overall verdict was \ACCEPT{} in five independent runs out of five, where it
had been \ACCEPT{} in three of five under the original all-seven scope. This
is a narrower claim than v1's, not a stronger one restated: RMS's
retained-architecture speed advantage is now stated to hold for $k>1$
against the same-work \texttt{eigh\_tridiagonal} peer, and no claim is made
about $k=1$ in either direction.

\section{Adversarial evaluation}
\label{sec:adversarial}

The seven declared instances share a shape: smooth potential, single well or
shallow double well, constant off-diagonal $-1/(2h^{2})$, moderate conditioning.
That is exactly the regime in which \texttt{DSTEBZ}'s scaling and block
splitting \citep{demmel1995} are idle and the native kernel's bare pivot floor
of $10^{-300}$ costs nothing. A comparison confined to it cannot support any
claim about robustness. This section reports what happens outside it. Three
findings follow; two are losses.

\subsection{Matrices built to defeat the kernel}
\label{sec:advmatrix}

We take the standard hard cases from the symmetric-tridiagonal literature and
hand both solvers the identical matrix at the same declared half-width
$2\times10^{-10}$, judging correctness against dense LAPACK---a different
algorithm on the full matrix, not a second bisection.

\begin{table*}[t]
\centering\small
\caption{Adversarial Field A. Errors are against dense \texttt{eigvalsh}.
Ratio is \texttt{DSTEBZ} time over native time, so values below $1$ are losses
for the native kernel. The glued Wilkinson matrix joins eight copies of
$W^{+}_{21}$ with an off-diagonal of $10^{-14}$; the denormal case sets one
off-diagonal to the pivot floor itself.}
\label{tab:adversarial}
\begin{tabular}{lrrrrrrr}
\toprule
Matrix & $n$ & $k$ & native err. & \texttt{DSTEBZ} err. & native (ms) & \texttt{DSTEBZ} (ms) & ratio \\
\midrule
Wilkinson $W^{+}_{21}$        & 21   & 4    & $3.00\text{e}{-11}$ & $3.02\text{e}{-11}$ & 0.035 & 0.043 & 1.239\\
Glued Wilkinson $8\times21$   & 168  & 4    & $2.26\text{e}{-11}$ & $2.19\text{e}{-11}$ & 0.074 & 0.101 & 1.363\\
Toeplitz $[-1,2,-1]$          & 4000 & 4    & $5.18\text{e}{-11}$ & $4.44\text{e}{-11}$ & 1.254 & 2.442 & 1.947\\
Scaled $10^{\pm10}$ (abs.\ tol) & 2000 & 4  & $2.48\text{e}{-5}$  & $2.29\text{e}{-5}$  & 1.662 & 2.891 & 1.739\\
Denormal split ($e=10^{-300}$) & 1000 & 4   & $3.35\text{e}{-11}$ & $5.19\text{e}{-11}$ & 0.347 & 0.541 & 1.561\\
Zero-coupling cluster         & 600  & 4    & $3.48\text{e}{-11}$ & $6.73\text{e}{-11}$ & 0.190 & 0.192 & 1.007\\
\midrule
\textbf{Glued Wilkinson}      & 168  & 64   & $4.80\text{e}{-11}$ & $7.93\text{e}{-11}$ & 0.725 & 0.346 & \textbf{0.478}\\
\textbf{Glued Wilkinson}      & 168  & 168  & $4.92\text{e}{-11}$ & $9.19\text{e}{-11}$ & 2.019 & 0.659 & \textbf{0.326}\\
Wilkinson $W^{+}_{401}$       & 401  & 64   & $9.05\text{e}{-11}$ & $6.17\text{e}{-11}$ & 1.827 & 3.775 & 2.066\\
Toeplitz $[-1,2,-1]$          & 4000 & 256  & $5.78\text{e}{-11}$ & $7.54\text{e}{-11}$ & 74.05 & 147.7 & 1.994\\
Toeplitz $[-1,2,-1]$          & 4000 & 1024 & $5.82\text{e}{-11}$ & $7.14\text{e}{-11}$ & 337.9 & 661.5 & 1.958\\
\bottomrule
\end{tabular}
\end{table*}

\Cref{tab:adversarial} separates two questions that are easy to conflate.

\emph{Accuracy is not the problem.} On every matrix the native kernel meets the
declared half-width and matches \texttt{DSTEBZ} to within a factor of two,
including on the denormal-split and zero-coupling cases that target the pivot
floor directly. The conjecture that the bare floor would produce wrong counts on
negligible off-diagonals is not borne out here. (The badly scaled case shows
errors of $2\times10^{-5}$ for \emph{both} solvers, which is what an
\emph{absolute} tolerance means on a matrix with $\|T\|_{1}\sim10^{10}$; it
separates nothing and is reported only so the row is not omitted.)

\emph{Speed is the problem, and exactly where predicted.} At $k=4$ the native
kernel is still ahead everywhere, though by less than on the physics suite. But
on the glued Wilkinson matrix at $k=64$ the ratio falls to $0.478$, and with all
$168$ eigenvalues requested to $0.326$: the native kernel is \emph{three times
slower}. The mechanism is not subtle. \texttt{DSTEBZ} detects the $10^{-14}$
coupling, splits the matrix into eight independent $21\times21$ blocks, and pays
$8\times21$ per sweep; the native kernel has no splitting and traverses all
$168$ pivots for every bisection step of every requested eigenvalue. The loss
grows with $k$ because the missing optimisation is per-eigenvalue.

This is the clearest statement of the kernel's boundary we can make. Its
advantage comes from batching several brackets through one diagonal sweep, and
that advantage inverts as soon as the opponent can make the sweep itself
shorter. Block splitting is a known, implementable technique; adding it is the
obvious next piece of work, and until it exists the native kernel should not be
offered for matrices with negligible off-diagonals.

\subsection{A non-smooth potential}
\label{sec:advnonsmooth}

Richardson extrapolation assumes the expansion \eqref{eq:expansion}, which is
not justified when $V$ has a kink. We therefore ran $V(x) = |x|$, whose spectrum
is known exactly in terms of Airy zeros: odd states at $E = -a_{n}/2^{1/3}$ and
even states at $E = -a'_{n}/2^{1/3}$, giving
$(0.80861652,\,1.85575708,\,2.57809613,\,3.24460762)$ for the four lowest.

All three pipelines return correct values. RMS gives a maximum error of
$1.23\times10^{-10}$ against a tolerance of $2\times10^{-8}$ with
$B = 3.24\times10^{-10}$, status \ACCEPT{}; the shipped and retained SciPy arms
give $1.28\times10^{-10}$ and $5.98\times10^{-10}$. The cost, however, is
visible: RMS needs $22$ solves rather than the $7$ of the smooth four-level
cases, and $21.2$\,ms rather than $4.4$. The single kink does not break the
extrapolation---plausibly because the gate detects the slower convergence and
simply refines further---but it triples the work, and this is one non-smooth
point, not a singular potential.

\subsection{A silent failure: the barrier that looks like a tail}
\label{sec:advfailure}

The most consequential result in this paper is negative and concerns none of the
things we set out to test.

Consider the symmetric double well $V(x) = \lambda(x^{2}-a^{2})^{2}$. With
$\lambda = 1$, $a^{2} = 4$ everything is fine: the tunnelling doublets are split
by $1.59\times10^{-5}$, all three pipelines resolve them, and RMS returns a
maximum error of $8.16\times10^{-10}$ at \ACCEPT{}. Now raise the barrier to
$a^{2} = 9$. The doublet splitting falls to $8.3\times10^{-12}$---below the
bisection half-width---and the reference spectrum, computed independently by a
sinc discrete variable representation \citep{colbert1992} at $N = 6000$, is
\begin{equation*}
\small (4.2144398,\ 4.2144398,\ 12.5272184,\ 12.5272184).
\end{equation*}

All three pipelines instead return
\begin{equation*}
\small (4.2144398,\ 12.5272184,\ 20.6574438,\ 28.5922506)
\end{equation*}
with status \ACCEPT{} and a reported diagnostic bound of
$1.58\times10^{-9} \le \eps$. The maximum error is $16.07$: nine orders of
magnitude above the declared tolerance, released with a passing gate.

The cause is not the kernel and not the extrapolation. It is the window. The
planner locates the minimum near $x = -3$ and returns the truncation window
$(-6.43,\ +0.43)$. That window contains one well and excludes the other
entirely; its right boundary sits inside the central barrier, where
$V \approx 77.7$ against $E_{3} \approx 28.6$. The decay gate of
\cref{sec:gate} asks whether the boundary lies deep in a classically forbidden
region, and the answer is yes---so the gate passes. The window witness of
\cref{sec:retain} then widens by a margin of $2s$, which is still not enough to
reach the second well, so the values barely move and the window witness passes
too. Both gates certify a truncation that has deleted half of the operator's
support, and what is returned is the correct spectrum of a \emph{different
problem}: the single well, whose level spacing of $\approx 7.9$ is exactly what
the returned values show.

Three points follow, and we state them plainly because they bear on the
architecture's central selling point.

\begin{enumerate}
\item \textbf{The failure is in the raw-input planner, not the solver.} Every
numerical component behaved correctly on the operator it was given. The claim
that a planner can be trusted to derive its own truncation from $(V,\theta,k,
\eps)$ is falsified by this instance.

\item \textbf{Both gates are unsound in the same way, and the second does not
catch the first.} The decay gate tests a \emph{local} property of the boundary
point---is $V(b)$ far above $E$?---and a barrier between two wells satisfies it
exactly as a decaying tail does. The window witness inherits the defect, because
a slightly wider window that is still inside the same barrier produces almost no
shift. Separated gates protect against separated failure modes; these two share
one.

\item \textbf{The declared suite could not have caught it.} It contains one
double well, \texttt{factorized\_sextic\_ground}, but with $k = 1$, so only the
ground state of one well is ever requested and the missing well never matters.

\end{enumerate}

The repair is not difficult to state, though we have not implemented it. The
gate must establish that the classically forbidden region \emph{extends to
infinity} rather than merely beginning at the boundary: for instance by
requiring $V$ to be monotone increasing beyond $b$ over a distance in which the
WKB action of \cref{rem:anticonservative} accumulates past $D_{\min}$, or by a
global scan for further minima of $V$ below $E_{k-1}$ outside the window before
any window is accepted. Until such a check exists, \ACCEPT{} on a potential with
more than one well should not be believed, and the honest scope of the planner
is single-well potentials with monotone tails.

\section{Scope of novelty}
\label{sec:novelty}

The following are classical and are claimed by no one here: central finite
differences; the symmetric tridiagonal reduction; the Sturm sequence; bisection;
Richardson extrapolation. Automatic domain truncation and automatic mesh
selection are likewise present in the established Sturm--Liouville packages
\citep{bailey2001,pruess1993,ledoux2005,pryce1993}. The architectural claim is
the composition
\begin{equation}
\boxed{
\begin{aligned}
&\text{raw-input planning}\\
&\quad + \text{validated retained brackets}\\
&\quad + \text{batched Sturm traversal}\\
&\quad + \text{multilevel correction retention}\\
&\quad + \text{separated mesh/window gates}
\end{aligned}}
\end{equation}
Retention in particular must be positioned carefully, because its closest
relative is not in the Sturm--Liouville literature at all. Nested iteration, and
its unidirectional coarse-to-fine variant, the cascadic multigrid method
\citep{bornemann1996,trottenberg2001}, solve on a coarse grid and carry the
result forward as the initial guess on the next finer grid; extrapolation
cascadic variants combine that inheritance with Richardson extrapolation
outright, and cascadic schemes for eigenvalue problems have been developed
since. The coarse-to-fine reuse in \cref{sec:brackets,sec:mesh} is therefore not
new as an idea, and we do not claim it is.

Two things do appear to differ, and they are what the claim should rest on.
First, the inherited object here is an \emph{interval revalidated by a Sturm
count}, not an initial vector that merely accelerates an iteration: an inherited
bracket that is wrong is detected and discarded before it can affect the answer,
whereas a bad initial guess in a cascadic scheme degrades silently into extra
iterations or a wrong limit. Second, the retention of \cref{sec:retain} crosses
a change of \emph{domain} at fixed spacing rather than a change of mesh, which
has no multigrid analogue, since multigrid refines the discretisation of a fixed
domain. Whether either difference is enough to constitute a contribution is a
judgement we leave to the reader; what we can say is that a comparison against
cascadic eigenvalue solvers, which we have not run, is the experiment that
would settle it.

With that caveat, the formulation of the contribution is:

\begin{quote}
We propose and evaluate a Retained Multilevel Sturm architecture for the lowest
eigenvalues of a tridiagonal discretisation, in which brackets and mesh
corrections computed at one level are retained and revalidated at the next, and
in which mesh resolution and domain-truncation stability are certified by
separate witnesses.
\end{quote}

We note one further methodological point inherited from the artefact and not
re-measured here. Timing LAPACK's MRRR driver (\texttt{DSTEMR}) through SciPy's
f2py wrapper is invalid for this comparison, because the wrapper declares a
dense $(n,n)$ eigenvector array even when eigenvectors are not requested; a
timing obtained that way measures an allocation the algorithm does not need. The
artefact therefore calls \texttt{DSTEMR} through \texttt{ctypes} with
\texttt{jobz='N'}. Any future work reporting MRRR comparisons should state which
route it used.

\section{Threats to validity}
\label{sec:threats}

Four items that were blocking in earlier drafts are marked \textbf{repaired};
each repair cost the method something, and the cost is stated. The rest remain
open.

\paragraph{T1. Mismatched comparator tolerance --- \textbf{repaired}, at a cost.}
Diagnosed in \cref{sec:matched}, repaired in Run C. The kernel geometric mean
falls from $2.560$ to $1.908$. Any speed figure quoted from an unmatched run
should be withdrawn. \Cref{rem:hang} records that the naive form of the repair
does not terminate, which is itself a result about the coupling between a
self-validating pipeline's acceptance gate and its solver tolerance.

\paragraph{T2. Unpaired, under-sampled, single-launch timing ---
\textbf{repaired}, at a cost.} The shipped Field-A design timed the arms in
separate loops, resampled them independently, and drew 11 samples from one
process. Run D interleaves the arms per repeat, resamples paired indices, and
pools $150$ pairs from three launches (\cref{sec:paired}). The five clear wins
survive with intervals roughly four times tighter;
\texttt{factorized\_sextic\_ground} becomes a reproducible measured loss and the
speed verdict falls to \HOLD{}.

\paragraph{T3. Comparator denied the idea under test ---
\textbf{repaired}, at the largest cost.} The shipped end-to-end comparator did
not implement retention, so its 12 solves were being reported as a property of
SciPy rather than of the comparator's author (\cref{sec:thirdarm}). With
retention ported in, the end-to-end geometric mean falls from $3.279$ to
$2.411$. This is the largest single correction in the paper.

\paragraph{T4. Fairness gate driven by a supplementary comparator ---
\textbf{repaired in code (v2), verdict narrowed rather than rescued}.}
Traced in \cref{sec:fairness} to a single ARPACK non-convergence on the
Morse case; v1 proposed a protocol-level repair (declare a primary
comparator) without implementing it. \Cref{sec:v2repair} (added in v2)
reports a different, implemented repair --- ARPACK switched to shift-invert
so it converges on all seven cases, verified against LAPACK --- which closes
the fairness gate deterministically. As v1 anticipated, this alone does not
rescue the overall verdict: the sole remaining cause is
\texttt{factorized\_sextic\_ground}'s speed, now attributed specifically to
$k=1$ and resolved by narrowing the verdict's scope to $k>1$, not by
overturning the $k=1$ result.

\paragraph{T5. Novelty of retention is not established.}
\Cref{sec:novelty} positions bracket inheritance against nested iteration and
cascadic multigrid \citep{bornemann1996,trottenberg2001} and identifies two
candidate differences, but we have run no comparison against a cascadic
eigenvalue solver. Until that exists, the architectural claim is a positioning
statement, not a demonstrated advance.

\paragraph{T6. ``Raw input, no tuning'' overstates the interface.}
The planner is free of \emph{user} tuning, not of tuning. The paths described
here contain at least fifteen fixed numeric constants---the $1025$-point probe,
the initial radius, the $8s$ half-width with its floor and cap, the $0.05\eps$
decay threshold, the $\kappa^{-1/4}$ scale, the factor $8$ and floor $64$ in
\eqref{eq:m0}, the $0.45\eps$ mesh gate, the $0.01\eps$ working width and its
$2\times10^{-12}$ floor, the $2s$ and $0.1$-width margins, the bisection
iteration cap, and the $10^{-300}$ pivot floor. Each was fixed once by the
author; none has a published sensitivity analysis. A referee is entitled to ask
what happens at $0.9\eps$ instead of $0.45\eps$, and we cannot answer.

\paragraph{T7. Comparator selection.}
The primary peer is \texttt{DSTEBZ}. The state of the art for requested-only
tridiagonal eigenvalues is MRRR \citep{dhillon2004}, and we report no MRRR
timings---only the methodological note of \cref{sec:novelty} about why the f2py
route cannot be timed. Beating the second-best routine while declining to
measure against the best is a legitimate objection, and the honest answer is
that the MRRR arm has not been run here.

\paragraph{T8. The declared suite is benign --- \textbf{addressed, with two
losses}.} All seven declared potentials are smooth and well conditioned.
\Cref{sec:adversarial} evaluates outside that regime and finds the boundaries:
the kernel loses by up to $3.1\times$ on block-splittable matrices
(\cref{sec:advmatrix}), a single kink triples the solve count
(\cref{sec:advnonsmooth}), and the planner returns a confidently \ACCEPT{}ed
answer that is wrong by $16.07$ on a high-barrier double well
(\cref{sec:advfailure}). Still untested: Coulomb and other singular potentials,
discontinuous $V$, and matrices requiring the scaling half of
\texttt{DSTEBZ}'s hardening \citep{demmel1995}.

\paragraph{T9. The boundary gate is unsound, not merely uncertified ---
\textbf{demonstrated}.} \Cref{rem:anticonservative} showed the decay score is
anti-conservative; \cref{sec:advfailure} shows it is worse than that. Because
the score tests a local property of the boundary point, a barrier separating two
wells satisfies it exactly as a decaying tail does, and the window witness
inherits the same blind spot. This is a soundness failure with a live
counterexample, not a missing proof. Any use of the word ``certified'' is
excluded until the gate establishes that the forbidden region extends to
infinity; the repair is sketched in \cref{sec:advfailure} and is not
implemented here.

\paragraph{T10. Richardson extrapolation is conditional; the bound is an
estimator.} Expansion \eqref{eq:expansion} is reasonable for smooth potentials.
\Cref{tab:bound} shows $B_{j}$ exceeds the audited error on all seven instances,
but by a stable factor, which identifies it as an error estimator with an
empirical safety margin rather than a bound (\cref{sec:boundsound}). Neither
result extends to non-smooth or singular potentials or to near-degenerate
clusters.

\paragraph{T11. Reproducibility record --- \textbf{improved, still incomplete}.}
\Cref{tab:host} records the replication host in full and Run D uses three
process launches. Two gaps remain: the recorded run's own \texttt{processor}
field is \texttt{unknown}, so Run A cannot be attributed to specific hardware;
and the replication host has one vCPU (\cref{tab:host}), so a regression that
would appear only when a multithreaded LAPACK is available to the comparator
cannot be detected here. Runs C and D should be repeated on a multi-core host.

\section{Conclusion}
\label{sec:conclusion}

Starting from raw problem data alone, the Retained Multilevel Sturm architecture
computed the declared low-lying spectra of seven one-dimensional Schr\"odinger
operators within their declared tolerances, with the reference values withheld
from every planning decision. On the harmonic instance the maximum audited
deviation is $1.465\times10^{-10}$ against a tolerance of $2\times10^{-8}$,
released under a diagnostic bound of $2.540\times10^{-9}$; across all seven the
released bound exceeds the audited error by factors of $7.7$ to $21.4$, which
identifies it as a well-behaved estimator rather than a proved bound.

The larger part of this paper is the correction of its own benchmark, and the
corrections were not free. Holding the comparator to the same bisection
half-width reduces the kernel-only geometric mean from $2.560$ to $1.908$.
Replacing the unpaired eleven-sample single-launch design with interleaved arms,
a paired bootstrap and $150$ pairs across three launches reduces it further to
$1.867$ and, more importantly, converts one instance into a reproducible
\emph{loss} to \texttt{DSTEBZ}, so the project's own speed verdict falls from
\ACCEPT{} to \HOLD{}. (\Cref{sec:v2repair}, added in v2, traces this
specific instance to \texttt{factorized\_sextic\_ground}'s $k=1$ and reports
that under a scope corrected to $k>1$ the verdict is \ACCEPT{} in five
independent runs out of five; the $k=1$ result itself is not overturned,
only excluded from a claim it cannot support.) Porting retention into the end-to-end comparator---giving
it the idea the architecture is proposing---reduces the end-to-end geometric
mean from $3.279$ to $2.411$. Two attempts to dismiss the remaining advantage
fail: it does not decay with problem size up to $n \approx 8\times10^{5}$, and
the diagnostic gate is never violated on the suite.

Evaluated outside the declared suite the boundaries appear quickly. The kernel's
advantage inverts on matrices its opponent can split into blocks---$3.1\times$
slower on a glued Wilkinson matrix with all eigenvalues requested---because its
advantage comes from batching brackets through one long sweep and block
splitting makes the sweep short. And on a symmetric double well whose central
barrier rises above the requested levels, the raw-input planner returns the
spectrum of a single well, with a passing decay gate, a passing window witness,
status \ACCEPT{}, and an error of $16.07$ against a tolerance of
$2\times10^{-8}$. The gate asks whether the boundary lies in a classically
forbidden region; a barrier answers yes exactly as a tail does. That is a
soundness failure in the component the architecture is named for, found by the
first potential we built to look for one, and it was invisible to a declared
suite whose only double well requests a single level.

What survives is narrower than what was originally claimed and, we think, more
\begin{equation}
\boxed{
\begin{aligned}
&\text{correct computation on seven smooth spectra}\\
&\quad + \text{a fair end-to-end advantage of }2.4\times\\
&\quad + \text{a kernel that loses where blocks split}\\
&\quad + \text{a truncation gate that is unsound}\\
&\quad + \text{a novelty claim not yet established}
\end{aligned}}
\end{equation}
and not a universal claim of superiority over SciPy. The honest scope of the
planner, until the gate of \cref{sec:advfailure} is repaired, is single-well
potentials with monotone tails.

\appendix

\section{Pipeline in pseudocode}
\label{app:pseudocode}

\begin{algorithm}[H]
\small
\caption{Retained Multilevel Sturm}
\label{alg:rms}
\begin{algorithmic}[1]
\Require $(V,\theta,k,\eps)$
\State $x_{\ast} \gets$ minimiser of $V$ on the probe grid
\State $\kappa \gets$ second central difference of $V$ at $x_{\ast}$;
       $s \gets \kappa^{-1/4}$
\State $w_{0} \gets \min(\max(8s,2),16)$;
       $[a_{0},b_{0}] \gets [x_{\ast}-w_{0},\,x_{\ast}+w_{0}]$
\State \textbf{if} decay score $D < D_{\min}$ \textbf{then} widen
       $[a_{0},b_{0}]$ and repeat
\State $M \gets M_{0}$; brackets $\gets$ Gershgorin bounds
\Repeat
  \State assemble $H_{M}$ from \eqref{eq:entries}
  \State \textbf{for} each active $j$: validate inherited bracket by Sturm
         count, doubling $r_{j}$ or falling back to global bounds
  \State batched bisection on \eqref{eq:pivot}--\eqref{eq:count} to width
         $\tau_{E}$
  \State $M \gets 2M$; inherit brackets
\Until{three meshes available and
       $\max_{j}\Delta_{\mathrm{mesh},j} \le 0.45\eps$}
\State $m \gets \max(2s,\,0.1(b_{0}-a_{0}))$;
       $[a_{1},b_{1}] \gets [a_{0}-m,\,b_{0}+m]$
\State $M_{\mathrm{expanded}} \gets M\,(b_{1}-a_{1})/(b_{0}-a_{0})$;
       one raw solve
\State $\Delta_{\mathrm{window},j}$ by \eqref{eq:windowwitness};
       $E_{j}$ by \eqref{eq:corrected}
\State \Return $E_{j}$ with status \ACCEPT{} iff
       $\max_{j}(\Delta_{\mathrm{mesh},j}+\Delta_{\mathrm{window},j})\le\eps$
\end{algorithmic}
\end{algorithm}

\section{Recorded configuration}
\label{app:config}

\begin{table*}[t]
\centering\small
\caption{Declared constants of the recorded run.}
\label{tab:config}
\begin{tabular}{ll}
\toprule
Quantity & Value \\
\midrule
Tolerance $\eps$ & $2\times10^{-8}$ \\
Levels $k$ & $4$ \\
Probe grid & $1025$ points, initial radius $8$ \\
Initial window & $[-8,8]$ \\
Verification window & $[-10,10]$ \\
Mesh ladder (initial window) & $128,256,512,1024,2048,4096$ \\
Expanded solve & $M = 5120$ (dimension $5119$) \\
Bisection width $\tau_{E}$ & $2\times10^{-10}$ \\
Mesh gate & $0.45\eps$ \\
Pivot floor & $10^{-300}$ \\
Kernel benchmark size & $M = 768$, $n = 767$, $k = 4$ \\
Kernel / end-to-end samples & $11$ / $21$ \\
Confidence intervals & percentile bootstrap, 95\% \\
\bottomrule
\end{tabular}
\end{table*}

\section{Replication host}
\label{app:host}

\begin{table*}[t]
\centering\small
\caption{Replication host for Runs B and C. Thread variables are pinned by the
harness before NumPy, SciPy, or Numba import, so the comparison measures
algorithms rather than the host thread scheduler.}
\label{tab:host}
\begin{tabular}{ll}
\toprule
Item & Value \\
\midrule
CPU & Intel Xeon @ 2.10\,GHz, 1 vCPU, 1 thread/core \\
RAM & 3\,GiB \\
Kernel & Linux 6.18.5, glibc 2.39 \\
Python & 3.12.3 \\
NumPy / SciPy & 2.4.4 / 1.17.1 \\
Numba / llvmlite & 0.66.0 / 0.48.0 \\
BLAS/LAPACK & OpenBLAS, 1 thread \\
Native kernel & compiled (\texttt{native\_kernel\_compiled = true}) \\
JAX & not installed (JAX comparator skipped) \\
Thread pinning & all BLAS/Numba thread vars $= 1$ \\
Repeats (Runs B, C) & 21 end-to-end, 11 kernel \\
Repeats (Run D) & 15 e2e, 50 kernel, $\times 3$ launches \\
Bootstrap & 4000 resamples, seed 20260727 \\
\bottomrule
\end{tabular}
\end{table*}

\begin{table*}[t]
\centering\small
\caption{Recorded run (Run A) environment, as committed. The \texttt{processor}
field is unset, which is the gap recorded under T7.}
\label{tab:hostA}
\begin{tabular}{ll}
\toprule
Item & Value \\
\midrule
Python & 3.13.13 \\
NumPy / SciPy & 2.4.6 / 1.17.1 \\
Numba & 0.65.1 \\
JAX / jaxlib & 0.11.0 / 0.11.0 \\
Platform & Linux 7.0.0, glibc 2.43 \\
Processor & \texttt{unknown} \\
\bottomrule
\end{tabular}
\end{table*}

\section{The Run C change}
\label{app:patch}

Run C differs from Run B by the following, and by nothing else. In the
end-to-end comparator:

\begin{lstlisting}[language=Python]
eigh_tridiagonal(
    diagonal, off_diagonal,
    select="i", select_range=(0, problem.modes - 1),
    eigvals_only=True, check_finite=False,
    tol=max(problem.tolerance * 0.01, 2.0e-12),   # matched working half-width
    lapack_driver="stebz",
)
\end{lstlisting}

\noindent and in the kernel audit, both solvers are held to the same width:

\begin{lstlisting}[language=Python]
def _matched_width(problem) -> float:
    return max(problem.tolerance * 0.01, 2.0e-12)

# native:
native_eigvals_from_tridiagonal(diagonal, off_diagonal, k, _matched_width(problem))
# comparator:
sla.eigh_tridiagonal(diagonal, off_diagonal, select="i",
                     select_range=(0, k - 1), eigvals_only=True,
                     check_finite=False,
                     tol=_matched_width(problem), lapack_driver="stebz")
\end{lstlisting}

\noindent Passing \texttt{tol=problem.tolerance} instead---the apparently
natural choice---does not terminate; see \cref{rem:hang}.

Run D adds two further pieces. The paired harness times every arm once per
repeat in a shuffled order and bootstraps paired repeat indices:

\begin{lstlisting}[language=Python]
for _ in range(repeats):
    rng.shuffle(names)              # arms interleaved, seeded
    for name in names:
        t = perf_counter(); calls[name]()
        samples[name].append(perf_counter() - t)

idx  = rng.integers(0, n, size=n)   # PAIRED resample
boot = median(b[idx]) / median(a[idx])
\end{lstlisting}

\noindent and the retained comparator arm replaces the second full Richardson
ladder by one same-spacing solve plus transport of the correction already
purchased---i.e.\ \eqref{eq:corrected}, with the window search, decay gate,
mesh ladder and \texttt{DSTEBZ} calls left untouched:

\begin{lstlisting}[language=Python]
ratio    = (exp[1] - exp[0]) / (win[1] - win[0])
n_exp    = max(64, int(round(finest * ratio)))
raw_exp  = solve(problem, exp, n_exp)      # ONE solve
shift    = abs(raw_exp - raw_finest)
values   = richardson_fine + (raw_exp - raw_finest)
\end{lstlisting}

\section*{Reproducibility statement}
\addcontentsline{toc}{section}{Reproducibility statement}
Run A quantities are read from the artefact's committed results file. Runs B,
C and D were produced by executing that same committed code on the host of
\cref{tab:host}: Run C with the single change of \cref{app:patch}, Run D
additionally with the paired-timing harness and the retained comparator arm of
\cref{app:patch}. Every figure in \cref{tab:kernel,tab:e2e,tab:scaling,tab:bound}
comes from a run executed for this paper, not copied from a prior report. All timing
figures are host-specific; the outstanding reproducibility gaps---the recorded
run's unset processor field, the absence of multiple independent process
launches, and the single-core replication host---are recorded under T7 and T8 of
\cref{sec:threats}.

\emph{Added in v2:} the ARPACK shift-invert repair of \cref{sec:v2repair} is
commit \texttt{72d5eec}, with the $10^{-9}$--$10^{-13}$ LAPACK-agreement
figure and the ARPACK-timing-inflation caveat both from the independent
re-verification of commit \texttt{acf4b51} (both merged to \texttt{main} as
part of pull request \#109); the instruction-count instrument and the $k>1$
scope correction are commits \texttt{43c0ac4} and \texttt{3823245} on pull
request \#110, which at the time of this writing has passed two rounds of
independent review but has \textbf{not yet been merged to \texttt{main}} ---
stated here so this is not read as already part of the artefact's default
branch.

\section*{Data and code availability}
\addcontentsline{toc}{section}{Data and code availability}
\begin{sloppypar}
The engine, the comparator pipeline, the same-operator
audit, the bootstrap routines, and the recorded
results file are distributed as part of the open-source
\emph{Information Discrete Mathematics} artefact under
the MIT licence, at
\url{https://github.com/morrocwi/information-discrete-math} (commits and
pull requests cited in this paper are resolvable at that repository). The
single-line change defining Run C is given verbatim in
\cref{app:patch}, so the controlled experiment of \cref{sec:bench} can be
reconstructed from the published sources without further information. The
v2 additions of \cref{sec:v2repair} (the ARPACK repair, the
\texttt{k1\_discrete\_readout} module, and the corresponding tests) are in
the same artefact and licence, at the commits given above. This manuscript
and its prose documentation are CC-BY-4.0; a showcase repository collecting
the paper source (both this manuscript and v1), benchmark scripts, and a
cross-reference from every claim in this paper to the script and data file
that produced it is at
\url{https://github.com/morrocwi/retained-sturm}.
\end{sloppypar}

\section*{Use of AI assistance}
\addcontentsline{toc}{section}{Use of AI assistance}
AI-based writing assistance was used for language editing and for structuring
the manuscript. All mathematical derivations, numerical results, interpretation,
and the identification of the defects reported in \cref{sec:threats} are the
responsibility of the author, who has verified every equation and every reported
figure. The v2 material of \cref{sec:v2repair} was produced with AI-assisted
implementation and analysis (the ARPACK repair, the instruction-count
instrument, and the corrected retention-mechanism claim), checked at each
step against the actual implementation and independently reviewed twice
before inclusion here; responsibility for the final claims remains the
author's, per the same standard as the rest of this statement.

\section*{Competing interests}
\addcontentsline{toc}{section}{Competing interests}
The author declares no competing interests.

\section*{Acknowledgements}
\addcontentsline{toc}{section}{Acknowledgements}
The author thanks Walancha Supantarika.

\begin{thebibliography}{99}

\bibitem[Anderson et al.(1999)]{lapack1999}
E.~Anderson, Z.~Bai, C.~Bischof, S.~Blackford, J.~Demmel, J.~Dongarra,
J.~Du~Croz, A.~Greenbaum, S.~Hammarling, A.~McKenney, and D.~Sorensen.
\newblock {\em LAPACK Users' Guide}.
\newblock SIAM, Philadelphia, 3rd edition, 1999.

\bibitem[Bailey et al.(2001)]{bailey2001}
P.~B. Bailey, W.~N. Everitt, and A.~Zettl.
\newblock Algorithm 810: The SLEIGN2 Sturm--Liouville code.
\newblock {\em ACM Transactions on Mathematical Software}, 27(2):143--192, 2001.

\bibitem[Barth et al.(1967)]{barth1967}
W.~Barth, R.~S. Martin, and J.~H. Wilkinson.
\newblock Calculation of the eigenvalues of a symmetric tridiagonal matrix by
the method of bisection.
\newblock {\em Numerische Mathematik}, 9:386--393, 1967.

\bibitem[Bornemann and Deuflhard(1996)]{bornemann1996}
F.~A. Bornemann and P.~Deuflhard.
\newblock The cascadic multigrid method for elliptic problems.
\newblock {\em Numerische Mathematik}, 75(2):135--152, 1996.

\bibitem[Colbert and Miller(1992)]{colbert1992}
D.~T. Colbert and W.~H. Miller.
\newblock A novel discrete variable representation for quantum mechanical
reactive scattering via the $S$-matrix Kohn method.
\newblock {\em Journal of Chemical Physics}, 96(3):1982--1991, 1992.

\bibitem[Demmel et al.(1995)]{demmel1995}
J.~W. Demmel, I.~Dhillon, and H.~Ren.
\newblock On the correctness of some bisection-like parallel eigenvalue
algorithms in floating point arithmetic.
\newblock {\em Electronic Transactions on Numerical Analysis}, 3:116--149, 1995.
\newblock (LAPACK Working Note 115.)

\bibitem[Dhillon and Parlett(2004)]{dhillon2004}
I.~S. Dhillon and B.~N. Parlett.
\newblock Multiple representations to compute orthogonal eigenvectors of
symmetric tridiagonal matrices.
\newblock {\em Linear Algebra and its Applications}, 387:1--28, 2004.

\bibitem[Efron(1979)]{efron1979}
B.~Efron.
\newblock Bootstrap methods: another look at the jackknife.
\newblock {\em Annals of Statistics}, 7(1):1--26, 1979.

\bibitem[Lam et al.(2015)]{lam2015}
S.~K. Lam, A.~Pitrou, and S.~Seibert.
\newblock Numba: a LLVM-based Python JIT compiler.
\newblock In {\em Proceedings of the Second Workshop on the LLVM Compiler
Infrastructure in HPC (LLVM '15)}, pages 1--6. ACM, 2015.

\bibitem[Ledoux et al.(2005)]{ledoux2005}
V.~Ledoux, M.~Van~Daele, and G.~Vanden~Berghe.
\newblock MATSLISE: A MATLAB package for the numerical solution of
Sturm--Liouville and Schr\"odinger equations.
\newblock {\em ACM Transactions on Mathematical Software}, 31(4):532--554, 2005.

\bibitem[Parlett(1998)]{parlett1998}
B.~N. Parlett.
\newblock {\em The Symmetric Eigenvalue Problem}.
\newblock Classics in Applied Mathematics 20. SIAM, Philadelphia, 1998.

\bibitem[Pruess and Fulton(1993)]{pruess1993}
S.~Pruess and C.~T. Fulton.
\newblock Mathematical software for Sturm--Liouville problems.
\newblock {\em ACM Transactions on Mathematical Software}, 19(3):360--376, 1993.

\bibitem[Pryce(1993)]{pryce1993}
J.~D. Pryce.
\newblock {\em Numerical Solution of Sturm--Liouville Problems}.
\newblock Oxford University Press, Oxford, 1993.

\bibitem[Richardson(1911)]{richardson1911}
L.~F. Richardson.
\newblock The approximate arithmetical solution by finite differences of
physical problems involving differential equations.
\newblock {\em Philosophical Transactions of the Royal Society A},
210:307--357, 1911.

\bibitem[Trottenberg et al.(2001)]{trottenberg2001}
U.~Trottenberg, C.~W. Oosterlee, and A.~Sch\"uller.
\newblock {\em Multigrid}.
\newblock Academic Press, London, 2001.

\bibitem[Virtanen et al.(2020)]{virtanen2020}
P.~Virtanen et al.
\newblock SciPy 1.0: fundamental algorithms for scientific computing in Python.
\newblock {\em Nature Methods}, 17:261--272, 2020.

\bibitem[Wilkinson(1965)]{wilkinson1965}
J.~H. Wilkinson.
\newblock {\em The Algebraic Eigenvalue Problem}.
\newblock Clarendon Press, Oxford, 1965.

\end{thebibliography}

\end{document}
